\cotefile{33}
\scanpage{2}
Nous prendrons comme prétexte la théorie des fonctions $L$ pour développer plus avant le formalisme de la cohomologie étale. Trois parties

a) Continuer l'étude \inserted{de la coh. étale} avec des coefficients de torsion, en développant en particulier, dans ce cadre, \deleted{\unsure{plus}} \deleted{\unsure{axiomatique}} \unsure{bien avant} le formalisme de dualité, \deleted{\unsure{donnant}} donnant des formules de Künneth et de Lefschetz.

b) Étendre le formalisme : des coefficients \uline{$l$-adiques}.

c) Appliquer les résultats au formalisme des fonctions $L$.

1°) \uline{Site étale} de $X$ ;
\[
\begin{array}{ll}
X_{\mathrm{ét}} : & \text{catégorie}\\
&\text{familles couvrantes}\\
&\text{faisceaux}\\
&\text{faisceaux abéliens}\\
&\text{cohomologie}
\end{array}
\]

2°) \uline{Exemple} $X-\unsure{\operatorname{Spec}k|\operatorname{Spm}k}$,
\[
X_{\mathrm{ét}}\simeq\widetilde X_{\mathrm{ét}}\simeq\operatorname{Ens}(\pi)
\qquad \pi=\operatorname{Gal}(\bar k/k)
\]

3°) $f:X\longrightarrow Y$ \unsure{définit}
\[
X_{\mathrm{ét}}-Y_{\mathrm{ét}}
\]
\[
\left\{\begin{aligned}
f^* &:Y_{\mathrm{ét}}\longrightarrow X_{\mathrm{ét}}\\
&:\widetilde Y_{\mathrm{ét}}\longrightarrow\widetilde X_{\mathrm{ét}}
\end{aligned}\right.
\]
\[
f_* :\widetilde X_{\mathrm{ét}}\longrightarrow\widetilde Y_{\mathrm{ét}}
\]
\margin{commute aux $\unsure{\varinjlim}$ \unsure{quelc.} et aux $\unsure{\varprojlim}$ finies}
\margin{commute aux $\unsure{\varprojlim}$ \unsure{quelc.}}
\mspage{1}
